% Gauss, Werke X/1 (1917), native113–114, printed106–107.
\SourceLeaf{113}\ArithmeticHead{106}{L}{ALGEBRA.\quad BRIEFWECHSEL.}
\NumberedItem{3}
\begin{center}\large\textsc{Gauss} an \textsc{Drobisch}.\end{center}
\begin{center}Wohlgeborner Herr\\Hochzuehrender Herr Professor.\end{center}
\noindent\hspace*{\parindent}Vor einigen Tagen ist mir das, durch Euer Wohlgeboren gütiges Schreiben
schon einige Wochen früher angemeldete Werk aus der \textsc{Dietrich}schen Buchhandlung zugekommen,
und ich eile nunmehro, Ihnen für beides meinen verbindlichsten Dank abzustatten.
Dieses schätzbare Werk wird gewiss dazu beitragen, in Deutschland die Liebe zu feineren
mathematischen Forschungen mehr zu verbreiten.

Es ist mir umso angenehmer gewesen, zu bemerken, dass Sie in die von mir vor einigen Jahren
zuerst in nuce bekannt gemachten[*)], obwohl schon seit fast 40 Jahren gehegten Grundansichten
über die imaginären Grössen eingegangen sind, da im Allgemeinen wenige für die Aufklärung der
Metaphysik der Mathematik Sinn zu haben scheinen, und namentlich die meisten sogenannten
Philosophen von Fach, wenn sie sich auf Mathematik einlassen, uns nur aegri somnia für Philosophie
verkaufen. Nur ist die Darstellung der imaginären Grössen in den Relationen der Puncte in plano
nicht sowohl ihr Wesen selbst, welches höher und allgemeiner aufgefasst werden muss, als vielmehr
das uns Menschen reinste oder vielleicht einzig ganz reine Beispiel ihrer Anwendung. Ich habe
öfters meine Theorie mündlich vorgetragen, und dann gefunden, dass sie sehr leicht aufgefasst
wird, und gar nichts abstruses behält.

Von \textsc{Fouriers} Werk habe ich in den hiesigen g[elehrten] A[nzeigen] 1833 Februar p. 321[**)]
eine Anzeige gegeben, indem ich darin dem was darin eigentlich verdienstlich ist Gerechtigkeit
widerfahren lassen, habe ich zugleich eine wenngleich nur leise Andeutung dessen was ich davon
tadle gegeben. Ich habe dort absichtlich mich über die captiöse Phrase, dass die fehlenden
Wurzeln imaginär werden, welche Phrase im Grunde ist, was die Engländer
\par\smallskip\noindent\rule{.25\linewidth}{.4pt}\par\smallskip
\begingroup\small
\noindent\hspace*{\parindent}[*) Anzeige der \textit{Theoria residuorum biquadraticorum,
Comm. secunda,} 1831, Werke II, S. 174 ff.]\par
\noindent\hspace*{\parindent}[**) Werke III, S. 110 ff.]\par
\endgroup
\clearpage
\SourceLeaf{114}\ArithmeticHead{107}{R}{GAUSS AN DROBISCH.}
\noindent unfair nennen, nicht weiter ausgelassen; aber ich kann auf das klarste beweisen,
dass ein Zusammenhang zwischen den \textls[120]{einzelnen} imaginären \textls[120]{Wurzelpaaren}\linebreak
mit \textls[120]{bestimmten} Stellen, wo Wurzeln ausfallen, \textls[120]{nicht} Statt findet,
und behalte mir vor, diess bei einer andern Gelegenheit zu entwickeln.

Es ist mir angenehm gewesen, dass Sie die Einfachheit und gleichsam Directheit meines ersten
Beweises die Wurzeln betreffend von 1799, anerkennen. Auch bin ich weit davon entfernt zu läugnen,
dass mein zweiter Beweis weitläufig ist. Mein Zweck dabei war hauptsächlich mit, nachzuweisen
was eigentlich den durchaus illusorischen Versuchen, mit denen selbst Geometer ersten Ranges
früher sich selbst getäuscht hatten, fehle, und wie es ergänzt werden müsse, und diess konnte
nicht ohne Ausführlichkeit geschehen. Hingegen kann ich den Vorwurf der Weitläufigkeit nicht
zugeben, wenn Sie ihn auch auf den Dritten erstrecken. Diesen halte ich für den kürzesten und
einfachsten von allen dreien. Auch ist derselbe im höchsten Grade sinnlich zu machen was ich
aber dort für ein hors d’oeuvre hätte halten müssen, wo ich alles in rein algebraischer Form
darstellen wollte. Übrigens aber finde ich den \textsc{Cauchy}schen Beweis in seiner Grundidee
eben so elegant, und möchte, wenn ich mich auf Einen beschränken sollte, selbst diesem den
Vorzug geben. Aber am lehrreichsten ist es wohl beide Grundideen zu entwickeln, wo in der That
die Betrachtung der gleichsam geometrischen Bedeutung beider für den Verstand etwas recht
ergötzliches hat. Eine concentrirte aber durchaus alles wesentliche enthaltende Darstellung
meines dritten Beweises hatte ich in den G. G. A. 1816 p. 338 f.[*)] gegeben.

Unter Bezeugung meiner ausgezeichneten Hochachtung und Ergebenheit habe ich die Ehre zu beharren
\par\medskip
\noindent\begin{minipage}[c]{.52\linewidth}Göttingen den 14. August 1834.\end{minipage}\hfill
\begin{minipage}[c]{.45\linewidth}
\begin{center}Ew. Wohlgeboren\\gehorsamer Diener\\C. F. \textsc{Gauss}.\end{center}
\end{minipage}\par\medskip
\noindent\hspace*{\parindent}H. Prof. \textsc{Möbius} bitte ich unter meiner besten Empfehlung um
die baldige Mittheilung der magnetischen Beobachtung am 6., 7. d. zu erinnern, die wir hier umso
begieriger erwarten, da hier am 7. ungemein auffallende Anomalien beobachtet sind.
\par\smallskip\noindent\rule{.25\linewidth}{.4pt}\par\smallskip
{\small\noindent\hspace*{\parindent}[*) Werke III, S. 107.]\par}
\vfill\noindent\hfill 14\textsuperscript{*}\par
\clearpage
